\section{Statistical Analysis Protocol}
\subsection{Units of analysis and dependence}
The model training run is the experimental unit.  Predictions on individual examples are
repeated observations produced by that unit and are not treated as independent model
replicates.  Likewise, the $\binom{k}{2}$ exact/exact pairs share models and therefore do
not constitute $\binom{k}{2}$ independent samples.  Primary uncertainty calculations
resample or aggregate at the model-seed level.  Example-level resampling may be used for a
conditional prediction analysis, but it must not replace uncertainty over training runs.

\subsection{Reference summaries}
For scalar distance $d$, the empirical exact/exact distribution is summarized by its
median, dispersion, and prespecified upper quantile.  Approximate/exact distances are
computed against every valid exact replica.  The principal descriptive statistic is the
difference between the approximate/exact median and the exact/exact median; a standardized
effect and empirical percentile provide scale-relative views.  When sample size permits,
bootstrap intervals resample exact models rather than treating dependent pairs as IID.

\subsection{Prospective exact-seed precision rule}
The number of exact replicas is selected by a deterministic precision analysis rather than
by inspecting approximate-method rankings.  Candidate counts are $k\in\{5,10,15\}$.  For
each primary metric, 5,000 repetitions sample whole exact-model seeds without replacement,
recompute the induced pairwise-distance mean, and form a 95\% interval.  A candidate count
passes when the interval half-width divided by its median is at most 0.20.  The selected
count is the first candidate passing for every primary metric.  The simulation seed is
73194. Ten exact seeds met the pilot rule for JS divergence in both calibrated full-class
conditions. This is a condition-specific pilot result, not a confirmatory or cross-dataset
seed-count guarantee.

The publication-scale profile uses a separate fail-closed decision.  Its exact-retraining
count cannot be smaller than 15, even when a pilot precision curve selects 10.  Original-model
and approximate-candidate replication are selected independently from
$\{5,8,10,15,20\}$ using validation-only, per-seed endpoint summaries.  Each count must satisfy
the 20\% relative interval-half-width rule in every frozen method--condition group.  Test
metrics, method rankings, winner labels, and confirmatory outputs are structurally forbidden
inputs.  At present, the exact-retraining floor is the only supported partial recommendation;
the other counts remain unset.  A hash-bound decision artifact must reproduce from its cited
evidence before any numeric count can enter the main profile.

\subsection{Equivalence and distinguishability}
Difference testing and equivalence testing answer different questions.  A failure to reject
equality does not establish equivalence.  For scalar endpoints with a scientifically
defensible margin $\Delta$, two one-sided tests or a compatible confidence-interval rule
will assess whether the entire interval lies in $[-\Delta,\Delta]$.  Margins must be frozen
before confirmatory outcomes are observed and reported in natural units.  If no defensible
margin exists, the analysis remains estimation-based and does not label models equivalent.

Calibration v1 proposed the condition-matched exact/exact 95th percentile as an upper margin
and zero as the lower bound.  Its prospectively fixed one-percent residual-class positive
control exceeded that margin in only 1/15 class-0 replicas and 0/15 class-5 replicas.  The
required rate was 80\%.  The v1 margins are therefore inadmissible.  This failure remains part
of the evidence record and cannot be repaired post hoc by widening the bound.

The unexecuted v2.1 design separates stochastic reference variability from the scientific
meaning of a discrepancy.  Let $R_{c,s}$ denote exact model seed $s$ in condition $c$, let
$\mathcal{R}_{c,-s}$ be its leave-one-out exact ensemble, and let $T_{\rho}$ swap the largest
and second-largest predicted probabilities on a deterministic class-stratified fraction
$\rho$ of a fixed validation probe.  For distance $d$, define the seed-level controlled excess
\begin{equation}
  \delta_{c,s,\rho} =
  \operatorname{median}_{r\in\mathcal{R}_{c,-s}} d(T_{\rho}(R_{c,s}),r)
  - \operatorname{median}_{r\in\mathcal{R}_{c,-s}} d(R_{c,s},r).
\end{equation}
The primary proposal uses $\rho=0.01$, which changes exactly 10 of 1,000 validation predictions
and, on the balanced Fashion-MNIST probe, exactly one prediction per class.  Fractions 0.005
and 0.02 are prespecified sensitivity anchors.  Fifteen calibration seeds estimate a candidate
margin as the one-sided 95\% lower bootstrap bound of the mean primary-control excess JS.
Ten disjoint holdout seeds must then yield at least 80\% primary effects above the margin, and
the holdout medians must increase strictly across the three severities.  All inputs are
validation-only; test predictions and approximate-method outputs are prohibited.

V2.1 is a design proposal, not a measured result or an approved margin.  The choice that one
percentage point of class-balanced top-1 disagreement is the smallest scientifically meaningful
effect requires an explicit human decision.  Until that review and a separate freeze amendment
exist, no equivalence test, equivalence label, or main-study launch is authorized.

\subsection{Multiplicity, exclusions, and failures}
No confirmatory family is currently active. The earlier proposed family containing JS
divergence and top-1 disagreement is suspended because disagreement was assay-insensitive
in the calibrated full-class conditions and the proposed equivalence margin failed. The
exploratory method pilot reports continuous estimates without multiplicity-adjusted
PASS/FAIL decisions. A later confirmatory family must be frozen prospectively; Holm or
another named correction will then be applied within that frozen family.

All planned seeds are retained unless a frozen quality rule excludes the run.
Infrastructure failures, numerical failures, and integrity violations are recorded
separately and are not statistical observations.  Exclusions must be determined without
reference to whether a run supports the hypothesis.  Missing runs are reported with their
configuration and remediation status.

\subsection{Exploratory versus confirmatory analysis}
P001 is retrospective and exploratory. Its seed-count convergence and identifiability
diagnostics can inform, but cannot validate, thresholds for C001. The two subsequent
calibration pilots are also non-confirmatory. The drafted full-class method pilot uses five
original seeds as the resampling units, ten exact references per condition, and 5,000
seed-level bootstrap repetitions. Exact/exact pairs are descriptive null values, not
independent observations. No equivalence or contradiction decision is permitted. Before
C001 begins, its hypotheses, endpoints, seed-generation procedure, correction family,
stopping rule, and failed-run policy must be timestamped. Any later change is an amendment,
and the affected result is labeled exploratory or post hoc.

\subsection{Sensitivity analyses}
Planned robustness checks vary the exact-envelope quantile, distance aggregation, seed
subsets, and calibration binning while preserving the primary analysis.  A conclusion that
depends on one arbitrary bin count or one exact replica is considered unstable.  For local
analyses, neighborhood definitions are frozen independently of the unlearning outcome to
avoid selecting regions that maximize observed discrepancy.
